\documentclass[11pt,a4paper]{article}

\usepackage{fontspec}
\usepackage[english]{babel}
\usepackage{amsmath, amsfonts, amssymb}
\usepackage{graphicx}
\usepackage{booktabs}
\usepackage{cite}
\usepackage{geometry}
\usepackage{xcolor}
\usepackage{url}
\usepackage{hyperref}

\geometry{margin=2.5cm}

\title{\textbf{Adaptive Sequenogram Extraction: Blind Deconvolution of Spectrograms via Physics-Informed Template-Based ConvNMF}}
\author{Your Name\thanks{Corresponding author: email@example.com} \\ 
        \small \textit{Department of Computational Signal Processing}}
\date{\today}

\begin{document}

\maketitle

\begin{abstract}
In complex acoustic environments, multipath propagation and reverberation significantly distort signal spectrograms, obscuring the temporal dynamics of sound emissions. We present a novel computational framework for extracting \textit{sequenograms}—high-resolution, sparse temporal activation vectors—from reverberant spectrograms using a dictionary of predefined templates. Our method employs a physics-informed Convolutional Non-Negative Matrix Factorization (ConvNMF) approach. Instead of using generic 2D kernels, we model the environmental echo using a frequency-dependent, causal 1D convolutional filter. This natively accounts for the physical phenomenon of frequency-dependent sound absorption and prevents the structural distortion of reference templates during optimization. Validated on synthetic datasets, the algorithm demonstrates robust performance in blind echo profile estimation and template localization. The proposed method provides a computationally efficient and adaptable tool for pattern discovery in large-scale acoustic monitoring.
\end{abstract}

\section{Introduction}
The precise localization and classification of signal patterns in time-frequency representations are fundamental challenges in various domains, including bioacoustics, industrial monitoring, and underwater signaling. In natural environments, such as caves, forests, or enclosed urban spaces, primary signals are inevitably followed by a series of reflections. These reflections form an ``echo tail'' that smears the acoustic energy across the time axis, effectively blurring the spectrogram. 

Traditional methods for signal detection, such as cross-correlation or standard image-based Convolutional Neural Networks (CNNs), are highly sensitive to such distortions. When the echo of one signal overlaps with the emission of a subsequent signal, standard detection algorithms often suffer from high false-positive rates and degraded timing accuracy.

To address this, we introduce the concept of a \textbf{sequenogram}: a 1D sparse activation vector that represents the exact emission times of specific templates, stripped of environmental artifacts. Our approach utilizes a physics-informed blind deconvolution technique to extract these sequenograms while simultaneously learning the adaptive echo profile of the recording environment.

\section{Materials and Methods}

\subsection{Physical and Generative Model}
Let $S \in \mathbb{R}_{\ge 0}^{F \times T}$ be the observed amplitude or power spectrogram, where $F$ is the number of frequency bins and $T$ is the number of time frames. We assume the generative process of the signal consists of two stages: emission and reverberation.

First, a ``clean'' (anechoic) spectrogram $V$ is generated. Given a set of $K$ known templates $W_k$, each corresponding to a specific signal class, the clean spectrogram is formed by convolving these templates with their respective non-negative activation vectors (sequenograms) $H_k \in \mathbb{R}_{\ge 0}^{T}$:
\begin{equation}
V(f, t) = \sum_{k=1}^{K} (H_k *_{t} W_k(f, \cdot))(t)
\end{equation}
where $*_{t}$ denotes convolution along the time axis.

Second, the clean spectrogram $V$ is subjected to environmental reverberation. Instead of using a generic 2D convolution that might erroneously shift energy across frequency bands, we apply a frequency-dependent 1D echo matrix $E \in \mathbb{R}_{\ge 0}^{F \times L_e}$. Each row of $E$ acts as an independent, causal impulse response for a specific frequency channel:
\begin{equation}
\hat{S}(f, t) = (V(f, \cdot) *_{t} E(f, \cdot))(t)
\end{equation}
This design natively reflects the physics of sound propagation: high-frequency sound waves attenuate more rapidly in the air than low-frequency waves.

\subsection{Physics-Informed Gauge Calibration of Loss Functionals}
To extract perfect sequenograms, we must impose structural constraints on the activation vectors $f(t) \in [0,1]$ (where $f(t)$ represents an arbitrary channel of $H$). Specifically, we require $f(t)$ to consist of isolated, binarized peaks (Dirac-like events) with an expected physical density $\rho_0$ (peaks per second). 

In standard multi-task learning, structural losses are combined using a linear combination $\mathcal{L}_{total} = \sum \alpha_i \mathcal{L}_i$, where $\alpha_i$ are manually tuned hyperparameters. This linear approach creates a fundamental geometric flaw: the gradients act as a flat plane, perpetually applying a restorative force even when a specific loss component reaches zero, potentially pushing the network into negative, unstable regimes.

Instead, we propose aggregating the loss components using a \textbf{Euclidean metric} $\mathcal{L}_{total} = \sqrt{\sum \mathcal{L}_i^2}$. The Euclidean geometry creates a conical potential well at the zero configuration. As a component $\mathcal{L}_i$ approaches zero, its partial derivative contribution naturally vanishes, dynamically shifting the optimization focus to the remaining unresolved constraints.

However, the Euclidean approach is only valid if all functional components $\mathcal{L}_i[f]$ reside in an isotropic space. We achieve this by formulating a \textit{Gauge Calibration} framework. Any valid functional must satisfy three strict physical conditions:
\begin{enumerate}
    \item \textbf{Zero Ground State:} $\mathcal{L}[f^*] = 0$ for the expected ideal configuration $f^*$.
    \item \textbf{Unit Variational Restoring Force:} The average magnitude of the variational derivative in the expected configuration must equal unity: $\langle | \frac{\delta \mathcal{L}}{\delta f(t)} | \rangle_{f^*} = 1$.
    \item \textbf{Discretization Invariance:} The functional must converge to a stable continuous integral as the sampling rate increases ($\Delta t \to 0$).
\end{enumerate}

Based on these principles, we derive two main penalty functionals for the sequenograms. Let the total physical time be $\tau_{max}$, and the expected event duration be $\tau_e$ (typically the limit of one discrete bin, $\tau_e = \Delta t$).

\subsubsection*{1. Sharpness Functional (Binarization)}
We penalize deviations from $0$ and $1$. The continuous functional is defined as:
\begin{equation}
\mathcal{L}_{sharp}[f] = \int_{0}^{\tau_{max}} f(t)\big(1 - f(t)\big) dt
\end{equation}
\textit{Proof of constraints:}
1) For $f(t) \in \{0, 1\}$, the integrand is strictly zero, satisfying the ground state condition. 
2) The variational derivative is $\frac{\delta \mathcal{L}_{sharp}}{\delta f(t)} = 1 - 2f(t)$. Since $f(t) \in \{0,1\}$, the norm $|1-2f(t)| = 1$ everywhere. Thus, its average is $1$.
3) Its discrete approximation $\mathcal{L}_{sharp} \approx \Delta t \sum_i f_i(1-f_i)$ explicitly incorporates $\Delta t$, ensuring invariance under varying sampling rates.

\subsubsection*{2. Isolation Functional (Window Penalty)}
To prevent peak blurring and excessively frequent emissions, we apply a local repulsion window. Let $w(t)$ be an indicator function equal to $1$ for $t \in [-\tau_w, \tau_w]$ and $0$ elsewhere, with a hollow center $w(0)=0$. The calibrated functional is:
\begin{equation}
\mathcal{L}_{iso}[f] = \lambda \int_{0}^{\tau_{max}} f(t) \big(w * f\big)(t) dt
\end{equation}
\textit{Proof of constraints:}
1) If peaks are separated by a distance greater than $\tau_w$, the convolution evaluates to $0$ at every peak location. Hence, $\mathcal{L}_{iso}[f^*] = 0$.
2) The variational derivative is $\frac{\delta \mathcal{L}_{iso}}{\delta f(t)} = 2\lambda (w * f)(t)$. In the expected configuration containing $N = \rho_0 \tau_{max}$ peaks of width $\tau_e$, the convolution equals $\tau_e$ inside the window of each peak. The non-zero domain covers $N \times 2\tau_w$ time. The average variational norm over $\tau_{max}$ is:
\begin{equation}
\langle | \frac{\delta \mathcal{L}_{iso}}{\delta f} | \rangle = \frac{1}{\tau_{max}} \left( N \cdot 2\tau_w \right) \cdot (2\lambda \tau_e) = 4 \lambda \rho_0 \tau_w \tau_e
\end{equation}
To enforce the unit variational restoring force (Condition 2), we uniquely determine the scaling constant $\lambda$:
\begin{equation}
\lambda = \frac{1}{4 \rho_0 \tau_w \tau_e}
\end{equation}
3) In discrete time, letting the window width be $W = \tau_w / \Delta t$ bins and peak width be $1$ bin ($\tau_e = \Delta t$), the discrete calibrated functional becomes:
\begin{equation}
\mathcal{L}_{iso} = \frac{\Delta t}{4 \rho_0 \tau_w} \sum_{i} f_i (w * f)_i
\end{equation}
This formulation preserves discretization invariance while adapting automatically to the physical parameters of the environment.

\subsubsection*{3. Unified Isotropic Objective}
Since both functionals are calibrated to project a unit gradient force around the solution, they form a perfectly isotropic error space. We combine them using the Euclidean norm:
\begin{equation}
\mathcal{L}_{reg} = \sqrt{\mathcal{L}_{sharp}^2 + \mathcal{L}_{iso}^2 + \epsilon}
\end{equation}
where $\epsilon$ is a small constant ensuring numerical stability at the origin. This regularization is added to the primary spectrogram reconstruction loss (KL-divergence).

\section{Results}

\subsection{Synthetic Validation (Proof of Concept)}
To verify the algorithm's capability to decouple overlapping signals and adaptive echoes, we constructed a synthetic dataset with known ground truth. The dataset included two templates occupying different frequency bands (low and high). Sparse activation spikes were generated for both templates, and a frequency-dependent exponential decay was applied as the echo. Finally, Gaussian noise was added to simulate real-world recording conditions.

The blind deconvolution model successfully reconstructed the sparse sequenograms with exact temporal precision, distinguishing overlapping calls. Furthermore, the learned matrix $E$ accurately matched the ground-truth decay rates, automatically discovering that higher frequencies decayed faster than lower ones. 

% [TODO: Insert Figure 1: Synthetic Spectrogram, Extracted Sequenograms, and Echo Profiles]

\subsection{Application to Real-World Acoustic Data}
% [TODO: Describe the results obtained from the 600,000x300 bat echolocation spectrogram. Mention how the templates were selected and how the extracted sequenograms improved signal timing accuracy compared to raw data.]

\section{Discussion}
The primary advantage of the proposed frequency-dependent 1D convolution over standard 2D deconvolution is its physical consistency. By prohibiting cross-frequency energy transfer within the echo model, we force the algorithm to explain temporal blurring strictly through reverberation. This prevents the model from maliciously altering the shape of the reference templates to minimize the loss function. 

Additionally, we introduced a Gauge Calibration framework for structural loss functions. By mathematically enforcing a \textit{unit variational restoring force} across all penalty terms, we mapped the optimization process into an isotropic space. This allowed us to aggregate structural constraints using Euclidean geometry ($\sqrt{\sum \mathcal{L}_i^2}$) rather than conventional linear combinations. We argue that if practitioners adopt these three physical rules (Zero Ground State, Unit Variational Force, and Discretization Invariance) for regularizing neural networks, the field could entirely eliminate the need for heuristic hyperparameter tuning when balancing multi-task objective functions.

\section{Conclusion}
The Adaptive Sequenogram Extraction framework offers a robust, physics-informed approach to pattern discovery in highly reverberant environments. By separating the intrinsic signal template from the external environmental echo, the method provides high-fidelity temporal activations. Its reliance on FFT-based convolutions and theoretically calibrated, scale-invariant loss functionals makes it computationally viable and hyperparameter-free for the large-scale analysis of continuous acoustic recordings.

\section*{Data and Code Availability}
The Python source code (utilizing PyTorch) and scripts for synthetic data generation are available in our public repository: \url{https://github.com/yourusername/AdaptiveSequenograms}.

\bibliographystyle{plain}
\bibliography{references}

\end{document}